% Page 005 — Margin notes + poetry

\seite{5}

\margmark{Die Bürger-\\meisterei\\Seffern.}%
mit welcher außerdem noch die Filialen Schleid u\ob
Heilenbach verbunden sind. Diese Orte bildeten zur\ob
Zeit der fr.\ Herrschaft eine eigene Bürgermeisterei unter\ob
dem Bürgermeister Joh.\ Weiler aus Seffern. Dieser\ob
Mann, er starb 1856, soll sich {[nicht]} wenig darauf zu gute ge\ohb
than haben, seiner Zeit in Schärpe und Federhut die\ob
neugebildete Bürgergarde drillen zu können.

\pend
\abschnitt{II. Alter}
\pstart

Wenn wir an die Frage über das Alter von Sefferweich\ob
herantreten, so dürfte man versucht sein, mit dem\ob
Dichter auszurufen: \enquote{Tot und stumm und öd und leer}.\ob
Keine Geschichte, keine Sage, kein monumentaler\ob
Bau giebt uns Aufschluß darüber; auch sonst ist noch\ob
niemand zurückgekommen, der uns Kunde brächte\ob
über den ersten Baumeister und die erste An\ohb
siedlung. Waren es Kelten, die hierherum gehaust\ob
haben oder waren es Germanen, Trevier oder Rö\ohb
mer, welche hier die ersten Fundamente gruben?\ob
Wer weiß es! Nicht einmal das Wort des Dichters\ob
findet hier seine Anwendung:

\pend
\begin{verse}
Nur ein halber Tempelriese\\
Trägt noch seine Quaderbalken,\\
Um die Giebel steigen Falken,\\
Epheu rankt sich um die Säule\\
Und die Natter und die Eidechs'
\end{verse}
\pstart
\pend
